%%% Task Manager — Yazılım Kalite Güvencesi
%%% Aşama 3: Kod Güncelleme — Güncelleme Raporu
%%% IEEE Conference / IEEEtran formatı
%%% Overleaf'te "New Project -> Blank Project" -> bu dosyayı yapıştır ve derle.
%%% Derleyici: pdfLaTeX (Overleaf -> Menu -> Compiler: pdfLaTeX).

\documentclass[conference]{IEEEtran}
\IEEEoverridecommandlockouts

% --- Türkçe karakter ve font desteği ----------------------------------------
\usepackage[utf8]{inputenc}
\usepackage[T1]{fontenc}
\usepackage[turkish]{babel}
\usepackage{lmodern}

% --- Tablo, listeleme ve dizgi paketleri ------------------------------------
\usepackage{cite}
\usepackage{amsmath,amssymb,amsfonts}
\usepackage{graphicx}
\usepackage{textcomp}
\usepackage{xcolor}
\usepackage{booktabs}
\usepackage{array}
\usepackage{hyperref}
\usepackage{url}
\usepackage{caption}

\hypersetup{
  colorlinks=true,
  linkcolor=black,
  urlcolor=blue,
  citecolor=black
}

\def\BibTeX{{\rm B\kern-.05em{\sc i\kern-.025em b}\kern-.08em
    T\kern-.1667em\lower.7ex\hbox{E}\kern-.125emX}}

\begin{document}

% ----------------------------------------------------------------------------
\title{Task Manager Uygulaması Üzerinde\\
Kod Güncelleme Raporu (Aşama 3)%
\thanks{Yazılım Kalite Güvencesi Dersi Dönem Projesi Aşama~3 Teslimi.}}

\author{%
\IEEEauthorblockN{Furkan Öztürk}
\IEEEauthorblockA{\textit{Yazılım Mühendisliği} \\
\textit{230229083}}
\and
\IEEEauthorblockN{Taha Yasin Çiçek}
\IEEEauthorblockA{\textit{Yazılım Mühendisliği} \\
\textit{230229088}}
\and
\IEEEauthorblockN{Betül Fikriye Turan}
\IEEEauthorblockA{\textit{Yazılım Mühendisliği} \\
\textit{220229073}}
\and
\IEEEauthorblockN{Elmas Güneş}
\IEEEauthorblockA{\textit{Yazılım Mühendisliği} \\
\textit{220229054}}
}

\maketitle

% ----------------------------------------------------------------------------
\begin{abstract}
Bu çalışma, Yazılım Kalite Güvencesi dönem projesinin üçüncü
aşaması (\emph{Kod Güncelleme}) kapsamında \emph{Task Manager} web
uygulamasında yapılan kalite iyileştirmelerini sunmaktadır. Aşama~2'de
\textbf{SonarQube~10.6 Community} ve \textbf{ESLint~8.57} ile elde edilen
toplam \textbf{124 bulgunun} (41 ESLint + 83 SonarQube) tamamı kod üzerinde
kapatılmıştır. Yalnızca \textbf{1~adet suppression} (yönerge limiti
3'tür) HTML5 native drag-and-drop API'sinin gerektirdiği olay dinleyicileri
için, gerekçesi belgelenerek kullanılmıştır. Sayısal sonuçlar:
\texttt{bugs}~14$\to$0, \texttt{code\_smells}~69$\to$0,
\texttt{security\_hotspots}~2$\to$0,
\texttt{reliability\_rating}~D$\to$A,
\texttt{sqale\_index}~348~dk$\to$0~dk. Aşama~1'de yazılan 98~test
(66~birim + 32~smoke) düzeltmeler sonrası tekrar koşturulmuş,
\textbf{regresyon gözlenmemiştir}. Aşama~4 için Aşama~2 sonunda konulan
tüm metrik hedefleri Aşama~3 sonunda aşılmıştır.
\end{abstract}

\begin{IEEEkeywords}
yazılım kalite güvencesi, kod güncelleme, SonarQube, ESLint, teknik borç,
refaktör, regresyon testi, TypeScript, React
\end{IEEEkeywords}

% ----------------------------------------------------------------------------
\section{Giriş}

Bu rapor, dönem projesinin dört aşamalı planının üçüncü aşamasını
belgelemektedir~\cite{b1,b2}. Aşama~2 İlk Statik Analiz Raporu'nda
SonarQube ve ESLint ile tespit edilen 124~bulgu Aşama~3 boyunca koda
uygulanan düzeltmelerle ele alınmış; istisnalar yönergeye uygun şekilde
gerekçelendirilmiştir. Düzeltmeler sonrasında ESLint ve SonarQube
yeniden çalıştırılmış, Aşama~1'de yazılan birim ve smoke testleri
tekrar koşturularak regresyon olmadığı doğrulanmıştır. Aşama~4'te
yapılacak son statik analiz, Aşama~3 ile aynı kod tabanı üzerinde fakat
\texttt{stage2} etiketiyle yeniden üretilecektir.

% ----------------------------------------------------------------------------
\section{Yöntem ve Çalışma Akışı}

Aşama~2 ile birebir aynı konfigürasyon dosyaları
(\texttt{.eslintrc.cjs}, \texttt{sonar-project.properties}) kullanılmış,
sürümler değiştirilmemiştir: ESLint~8.57.0, SonarQube~10.6 Community,
sonar-scanner~3.1.0, Node.js~22.15.0, Java~21. Çalışma akışı şu
sıradadır: (i)~bulgu listesi öncelik sırasına göre alınmış, (ii)~her
bulgu kategorisi için kod üzerinde düzeltme yapılmış, (iii)~her büyük
düzeltmeden sonra \texttt{npm test} ile regresyon kontrol edilmiştir,
(iv)~sonunda \texttt{npm run lint:stage2} ve
\texttt{sonar-scanner} yeniden çalıştırılmış, sonuçlar
\texttt{reports/stage2/} altında saklanmıştır.

% ----------------------------------------------------------------------------
\section{Özet Sonuçlar}
\label{sec:ozet}

Aşama~2 ve Aşama~3 arasındaki sayısal değişim Tablo~\ref{tab:summary}'de
sunulmuştur.

\begin{table}[htbp]
\caption{Aşama 2 / Aşama 3 metrik karşılaştırması}
\label{tab:summary}
\centering
\begin{tabular}{@{}lrrr@{}}
\toprule
\textbf{Metrik} & \textbf{Aşama 2} & \textbf{Aşama 3} & \textbf{$\Delta$} \\
\midrule
ESLint error                      &   15 &   \textbf{0} & $-15$ \\
ESLint warning                    &   26 &   \textbf{0} & $-26$ \\
ESLint toplam                     &   41 &   \textbf{0} & $-41$ \\
SonarQube BUG                     &   14 &   \textbf{0} & $-14$ \\
SonarQube CODE\_SMELL             &   69 &   \textbf{0} & $-69$ \\
\texttt{vulnerabilities}          &    0 &   0 & 0 \\
\texttt{security\_hotspots}       &    2 &   \textbf{0} & $-2$ \\
SonarQube açık issue              &   83 &   \textbf{0} & $-83$ \\
\texttt{reliability\_rating}      & D (4.0) & \textbf{A (1.0)} & $\uparrow$ \\
\texttt{sqale\_rating}            & A & A & korundu \\
\texttt{security\_rating}         & A & A & korundu \\
\texttt{sqale\_index}             & 348 dk & \textbf{0 dk} & $-348$ \\
\texttt{sqale\_debt\_ratio}       & \%0.2 & \textbf{\%0.0} & $-0.2$ \\
\texttt{duplicated\_lines\_density} & \%1.3 & \%1.0 & $-0.3$ \\
\texttt{ncloc}                    & 5153 & 5308 & $+155$ \\
\bottomrule
\end{tabular}
\end{table}

Toplam 124~bulgudan 124'ü kapatılmış olup yönergenin izin verdiği üç
suppression hakkından yalnızca biri kullanılmıştır
(Bölüm~\ref{sec:suppress}).

% ----------------------------------------------------------------------------
\section{Çözülen Bulguların Listesi}
\label{sec:cozumler}

Bu bölüm, kapatılan bulguları kategori bazında sunar. Her kategori için
\emph{problem özeti}, \emph{yapılan değişiklik} ve \emph{nasıl
doğrulandığı} (test, manuel kontrol, yeniden tarama) belirtilmiştir.

\subsection{ESLint Bulguları (41 adet)}

\subsubsection{B-E1. \texttt{@typescript-eslint/no-explicit-any} (36 adet)}
\textbf{Problem özeti:} \texttt{any} tipi kullanıldığı her yerde TypeScript
tip güvenliğini devre dışı bırakıyor; alan/metot çağrı hatalarını gizliyordu.

\textbf{Dağılım:}
\begin{itemize}
\item \texttt{src/services/projectService.ts} -- 19 adet
\item \texttt{client/src/pages/TasksPage.tsx} -- 11 adet
\item \texttt{client/src/pages/AdminPage.tsx} -- 5 adet
\item \texttt{src/services/taskService.ts} -- 3 adet
\item \texttt{src/routes/tasks.ts} -- 1 adet
\end{itemize}

\textbf{Yapılan değişiklik:} Domain tiplerinden uygun \texttt{interface},
\texttt{type-alias} veya \texttt{unknown}+narrowing kullanılarak
\texttt{any} kaldırıldı. Yeni tipler: \texttt{Id}, \texttt{ProjectMemberRow},
\texttt{SqliteError}, \texttt{TaskRow}, \texttt{TaskFormData}.

\textbf{Doğrulama:} \texttt{npm run lint:stage2} 0 hata/uyarı; backend ve
frontend için \texttt{tsc --noEmit} temiz; 98/98 test geçiyor.

\subsubsection{B-E2. \texttt{react/no-unescaped-entities} (2 adet)}
\textbf{Problem özeti:} JSX içinde literal \texttt{"} karakteri.
\textbf{Yapılan değişiklik:} \texttt{TasksPage.tsx} arama sonuç başlığında
\texttt{\&quot;} ile kaçırıldı.
\textbf{Doğrulama:} ESLint stage2 raporu temiz.

\subsubsection{B-E3. \texttt{@typescript-eslint/no-unused-vars} (2 adet)}
\textbf{Problem özeti:} \texttt{tests/unit/projectService.test.ts} ve
\texttt{tests/unit/taskService.test.ts} dosyalarında kullanılmayan import.
\textbf{Yapılan değişiklik:} Gereksiz importlar kaldırıldı.
\textbf{Doğrulama:} ESLint stage2 temiz; 66 birim testin tamamı geçmeye
devam ediyor.

\subsubsection{B-E4. \texttt{react-hooks/exhaustive-deps} (1 adet)}
\textbf{Problem özeti:} \texttt{TasksPage.tsx} içinde \texttt{useEffect}
bağımlılık dizisi \texttt{loadTasks} fonksiyonunu kapsamıyordu.
\textbf{Yapılan değişiklik:} \texttt{loadTasks} \texttt{useCallback}
ile sarmalandı; \texttt{useEffect} bağımlılık dizisi düzeltildi.
\textbf{Doğrulama:} ESLint kuralı tetiklenmiyor; tarayıcıda hızlı ekleme
ve filtre senaryoları manuel olarak doğrulandı.

\subsection{SonarQube Bulguları (83 adet)}

\subsubsection{B-S1. \texttt{S2871} -- CRITICAL BUG -- \texttt{.sort()}
karşılaştırıcısı eksik}
\textbf{Konum:} \texttt{src/routes/tasks.ts:130} (eski).
\textbf{Problem özeti:} JavaScript varsayılanı sayıları string olarak
sıralıyor; atanan ID'lerin karşılaştırması güvenilmez.
\texttt{reliability\_rating}'i D'ye çeken tek CRITICAL bulgu.
\textbf{Yapılan değişiklik:} \texttt{buildAssigneeIdsKey} yardımcı
fonksiyonunda \texttt{[\dots ids].sort((a,b) => a-b)} sayısal
karşılaştırıcı verildi.
\textbf{Doğrulama:} Yeniden tarama: \texttt{bugs=0},
\texttt{reliability\_rating=A}.

\subsubsection{B-S2. \texttt{S3776} -- CRITICAL CODE\_SMELL -- Cognitive
Complexity}
\textbf{Konum:} \texttt{src/routes/tasks.ts:113} (eski) PATCH \texttt{/:id}
handler. Cognitive complexity~23, eşik~15.
\textbf{Yapılan değişiklik:} Mantık iki saf yardımcıya bölündü:
\texttt{buildAssigneeIdsKey} ve \texttt{describePatchChanges}. Handler
artık 5~satırlık bir akışa sahip, eşik altında.
\textbf{Doğrulama:} Yeniden tarama: \texttt{S3776} raporda yok.
Smoke test~3 (\texttt{PATCH /tasks/:id}) ve birim testler yeşil.

\subsubsection{B-S3. \texttt{S6606} -- \texttt{??} yerine \texttt{||} (25 adet)}
\textbf{Problem özeti:} \texttt{||} operatörü \texttt{0}, \texttt{''} ve
\texttt{false} gibi geçerli \emph{falsy} değerleri de filtreliyor; ör.\ task
sayacı \texttt{0} iken yanlışlıkla varsayılan kullanılıyordu.
\textbf{Yapılan değişiklik:} 25 noktada \texttt{||} $\to$ \texttt{??}.
Etkilenen dosyalar: \texttt{AdminPage.tsx} (8), \texttt{TasksPage.tsx} (7),
\texttt{ReportsPage.tsx} (3), \texttt{server.ts} (2),
\texttt{routes/tasks.ts} (1), \texttt{labelService.ts} (1),
\texttt{projectService.ts} (1), \texttt{app.ts} (1),
\texttt{LoginPage.tsx} (1).
\textbf{Doğrulama:} Yeniden tarama: \texttt{S6606} raporda yok.

\subsubsection{B-S4. \texttt{S6853} -- Form etiketi--kontrol bağlanmamış (15 adet)}
\textbf{Problem özeti:} \texttt{<label>} öğelerinin \texttt{htmlFor}
özelliği yoktu; ekran okuyucular etiketi ilgili \texttt{<input>} ile
eşleştiremezdi.
\textbf{Yapılan değişiklik:} \texttt{LoginPage.tsx} (2),
\texttt{RegisterPage.tsx} (3) ve \texttt{TasksPage.tsx} (10)
formlarındaki tüm etiketlere benzersiz \texttt{id} ve \texttt{htmlFor}
verildi. Görsel olmayan etiketler için \texttt{visually-hidden} sınıfı
kullanıldı.
\textbf{Doğrulama:} Yeniden tarama: \texttt{S6853} raporda yok.

\subsubsection{B-S5. \texttt{S6759} -- \texttt{readonly} eksik prop (4 adet)}
\textbf{Yapılan değişiklik:} \texttt{StatusBadge}, \texttt{PriorityBadge},
\texttt{TaskFormProps}, \texttt{TaskDetailProps}, \texttt{LoginPage.Props},
\texttt{RegisterPage.Props} arabirimlerinde tüm alanlara \texttt{readonly}
eklendi.
\textbf{Doğrulama:} Yeniden tarama: \texttt{S6759} raporda yok.

\subsubsection{B-S6. \texttt{S6848} + \texttt{S1082} -- Etkileşimsiz öğe
üstünde click handler (27 adet)}
\textbf{Problem özeti:} \texttt{<div>} olarak yazılmış modal arkaplanları,
sidebar overlay, proje listesi öğeleri ve task kartları klavye/ekran
okuyucu kullanıcılarına ulaşılabilir değildi.
\textbf{Yapılan değişiklik:} Erişilebilirlik refaktörü:
\begin{itemize}
\item Modal overlay (4 adet): dış \texttt{<div>} salt-stil; arkaplan
  tıklamasını yakalayan ayrı \texttt{<button class="modal-overlay-backdrop">}
  (CSS \texttt{position:absolute; inset:0}) eklendi.
\item Sidebar overlay (mobil): \texttt{<div>} $\to$ \texttt{<button>}.
\item Proje sidebar öğeleri: dış \texttt{<div>} $\to$ \texttt{<button>};
  silme butonu kardeş öğe olarak ayrıldı.
\item Task kartı: \texttt{<div>} $\to$ \texttt{<article>}; içine
  ``stretched-link'' \texttt{<button class="task-card-stretch">}
  yerleştirildi (mutlak konum, kart yüzeyini kaplar). Düzenle/sil butonları
  \texttt{z-index:2} ile kalır.
\item Modal iç \texttt{<div>}'lerinden gereksiz
  \texttt{onClick=stopPropagation} kaldırıldı.
\end{itemize}
\textbf{Doğrulama:} Yeniden tarama: \texttt{S6848} ve \texttt{S1082}
raporda yok. Klavye akışı manuel olarak doğrulandı: Tab ile odak,
Enter ile aktivasyon çalışıyor.

\subsubsection{B-S7. \texttt{S3863} -- Yinelenen \texttt{import} (3 adet)}
\textbf{Yapılan değişiklik:} \texttt{src/routes/tasks.ts}'de iki
\texttt{activityService} importu birleştirildi;
\texttt{TasksPage.tsx}'te \texttt{lucide-react} ve \texttt{'../types'}
importları tek bloklara toplandı.

\subsubsection{B-S8. \texttt{S3358} -- İç içe ternary (2 adet)}
\textbf{Yapılan değişiklik:}
\texttt{src/routes/tasks.ts:31} (eski): \texttt{parseProjectIdFilter()};
\texttt{TasksPage.tsx:51} (eski): \texttt{statusBadgeClass()} yardımcı
fonksiyonlarına çıkarıldı.

\subsubsection{B-S9. \texttt{S3696} -- Literal \texttt{throw} (1 adet)}
\textbf{Konum:} \texttt{client/src/api.ts:11}.
\textbf{Yapılan değişiklik:} \texttt{class ApiError extends Error}
tanımlandı; \texttt{throw new ApiError(res.status, body)} kullanıldı.

\subsubsection{B-S10. \texttt{S3626} -- Gereksiz \texttt{return} (1 adet)}
\textbf{Yapılan değişiklik:} Arrow fonksiyonun sonundaki gereksiz
\texttt{return;} satırı kaldırıldı.

\subsubsection{B-S11. Security Hotspots: \texttt{S5852} ve \texttt{S5689}}
\textbf{S5852 -- regex super-linear backtracking}
(\texttt{src/services/authService.ts:30}): Eski regex tabanlı e-posta
doğrulayıcı kaldırıldı; yerine RFC 5321 sınırlarını (254/64) uygulayan
saf \texttt{isValidEmail()} fonksiyonu eklendi (lineer zaman, backtracking
yok).
\textbf{S5689 -- framework version disclosure}
(\texttt{src/app.ts:11}): \texttt{app.disable('x-powered-by')} eklenerek
\texttt{X-Powered-By} HTTP başlığı kapatıldı.
\textbf{Doğrulama:} Yeniden tarama: \texttt{security\_hotspots=0}.
\texttt{authService.test.ts}'in 10 testi yeniden geçti.

% ----------------------------------------------------------------------------
\section{Regresyon Test Kanıtı}
\label{sec:regresyon}

Aşama~1'de yazılan testlerin tamamı (66~birim + 32~smoke = \textbf{98 test})
Aşama~3 düzeltmeleri sonrası tekrar koşturulmuştur. Çıktı,
\texttt{asama\_3/reports/stage2/regression-tests.log} dosyasında
saklıdır. Özet sonuç Tablo~\ref{tab:tests}'de sunulmuştur.

\begin{table}[htbp]
\caption{Regresyon test sonucu (\texttt{npm test})}
\label{tab:tests}
\centering
\begin{tabular}{@{}lrr@{}}
\toprule
\textbf{Test paketi} & \textbf{Test} & \textbf{Sonuç} \\
\midrule
\texttt{authService.test.ts}    & 10 & \checkmark \\
\texttt{projectService.test.ts} &  7 & \checkmark \\
\texttt{taskService.test.ts}    & 12 & \checkmark \\
\texttt{labelService.test.ts}   &  9 & \checkmark \\
\texttt{commentService.test.ts} &  8 & \checkmark \\
\texttt{reportService.test.ts}  &  8 & \checkmark \\
\texttt{api.smoke.test.ts}      & 32 & \checkmark \\
\midrule
\textbf{Genel}                  & \textbf{98} & \textbf{98 / 98} \\
\bottomrule
\end{tabular}
\end{table}

\noindent Ham çıktı:
\begin{small}
\begin{verbatim}
Test Files  7 passed (7)
     Tests  98 passed (98)
  Duration  1.44s
\end{verbatim}
\end{small}

% ----------------------------------------------------------------------------
\section{Kullanılan Suppression / Ignore}
\label{sec:suppress}

Yönerge en fazla \textbf{3} suppression'a izin vermektedir;
\textbf{1~adet} kullanılmıştır.

\subsection{Suppression \#1}

\begin{itemize}
\item \textbf{Bulgu kimliği / kural adı:} \texttt{typescript:S6847}
  ``\emph{Non-interactive elements should not be assigned mouse or
  keyboard event listeners}''.
\item \textbf{Kapsam:} \texttt{client/src/pages/TasksPage.tsx}.
\item \textbf{Mekanizma:} \texttt{sonar-project.properties} içinde
  \texttt{sonar.issue.ignore.multicriteria.s6847dnd}.
\end{itemize}

\subsubsection{Neden istisna?}
HTML5 \emph{native} drag-and-drop API'si \emph{drop target}'ta
\texttt{onDragOver}~+~\texttt{onDrop}, \emph{drag source}'ta
\texttt{onDragStart} olaylarını doğrudan elementin üstünde gerektirir;
bu olaylar başka bir interaktif öğeye taşınamaz. SonarQube'un
\texttt{S6847} kuralı sürükle-bırak olaylarını da ``mouse listener''
sayıp etkileşimsiz \texttt{<section>} ve \texttt{<article>} öğelerinde
tetiklemektedir.

\subsubsection{Risk değerlendirmesi}
\textbf{Düşük.} Erişilebilirlik açısından risk yoktur:
\begin{itemize}
\item Klavye/ekran okuyucu kullanıcıları için detayı açma akışı task kart
  üstüne eklenen \emph{stretched-link} \texttt{<button>}
  (\texttt{.task-card-stretch}) ile sağlanmaktadır
  (Tab+Enter çalışıyor).
\item Durum değiştirmek için modal içindeki ``Durum'' \texttt{<select>}
  elemanı alternatif yol olarak mevcuttur.
\item Sürükle-bırak yalnızca fare/dokunmatik için ek kolaylıktır;
  klavye yolundan vazgeçirilmemiştir.
\end{itemize}

\subsubsection{Alternatif çözüm neden uygulanmadı?}
İki olası alternatif değerlendirilmiştir:
\begin{enumerate}
\item \emph{Task kartını ve kanban kolonunu \texttt{<button>} yapmak.}
  HTML standardı \texttt{<button>} içine başka \texttt{<button>}
  yerleştirmeye izin vermemektedir; oysa task kartının içinde Düzenle/Sil
  butonları bulunmaktadır. Ayrıca kolonun başlığı, sayaç ve görev listesi
  içermesi onu \texttt{<button>} olmaktan semantik olarak alıkoyar.
\item \emph{Native HTML5 \texttt{<dialog>} + dnd-kit / react-dnd
  kütüphaneye geçmek.} Yeni bağımlılık, yeni öğrenme eğrisi ve mevcut
  CSS'in baştan yazılması gerekir; üç haftalık Aşama~3 takvimi içinde
  gerekçeli görülmemiştir. Mevcut native DnD pattern'ı, \emph{stretched-link}
  butonlar üzerinden zaten erişilebilir alternatif sunmaktadır.
\end{enumerate}

% ----------------------------------------------------------------------------
\section{Aşama 4 Hedeflerinin Erken Karşılanması}
\label{sec:hedef}

Aşama~2 raporunda Aşama~4 için konulmuş hedefler ile Aşama~3 sonu
sonuçları Tablo~\ref{tab:hedefler}'de karşılaştırılmıştır. Tüm hedefler
Aşama~3'te erken karşılanmıştır.

\begin{table}[htbp]
\caption{Aşama 2 hedefleri vs.\ Aşama 3 gerçekleşen}
\label{tab:hedefler}
\centering
\begin{tabular}{@{}lrrl@{}}
\toprule
\textbf{Metrik} & \textbf{Aşama 2} & \textbf{Hedef} & \textbf{Aşama 3} \\
\midrule
ESLint error             & 15 & 0 & \textbf{0} \\
ESLint warning           & 26 & $\leq$ 5 & \textbf{0} \\
SonarQube BUG            & 14 & $\leq$ 2 & \textbf{0} \\
SonarQube CODE\_SMELL    & 69 & $\leq$ 20 & \textbf{0} \\
BLOCKER+CRITICAL CODE\_SMELL & 1 & 0 & \textbf{0} \\
\texttt{duplicated\_lines\_density} & \%1.3 & $\leq$ \%1.0 & \textbf{\%1.0} \\
\texttt{reliability\_rating} & D & A & \textbf{A} \\
\texttt{sqale\_rating}   & A & A & A \\
\texttt{security\_rating} & A & A & A \\
\texttt{sqale\_index}    & 348 dk & $\leq$ 120 dk & \textbf{0 dk} \\
\texttt{security\_hotspots} & 2 & 0 & \textbf{0} \\
\bottomrule
\end{tabular}
\end{table}

% ----------------------------------------------------------------------------
\section{Teslimat Paketi}

Aşağıdaki dosyalar Aşama~3 teslim paketini oluşturur:

\begin{itemize}
\item \texttt{asama\_3/GUNCELLEME\_RAPORU.md} -- Markdown sürümü.
\item \texttt{asama\_3/GUNCELLEME\_RAPORU.tex} -- bu LaTeX kaynağı.
\item \texttt{asama\_3/reports/stage2/eslint.json} -- ESLint çıktısı
  (0~bulgu).
\item \texttt{asama\_3/reports/stage2/sonarqube-issues.json} -- SonarQube
  çıktısı (0 OPEN issue).
\item \texttt{asama\_3/reports/stage2/sonarqube-measures.json} --
  13~SonarQube metriği.
\item \texttt{asama\_3/reports/stage2/regression-tests.log} -- 98/98
  test çıktısı.
\item Düzeltilmiş kod tabanı -- \texttt{src/}, \texttt{client/src/},
  \texttt{tests/} ve \texttt{sonar-project.properties}.
\end{itemize}

% ----------------------------------------------------------------------------
\section{Sonuç}

Aşama~2 raporundaki 124 bulgunun tamamı kod üzerinde kapatılmış;
yalnızca HTML5 native drag-and-drop için 1~suppression kullanılmıştır.
\texttt{reliability\_rating} D'den A'ya yükselmiş, teknik borç (\texttt{sqale\_index})
348~dakikadan 0'a inmiş, güvenlik incelemesi gerektiren noktalar gerçek
düzeltmelerle (regex backtracking ve framework version disclosure)
ortadan kaldırılmıştır. 98~birim ve smoke testinin tamamı Aşama~3
düzeltmeleri sonrası başarıyla geçmektedir; regresyon yoktur. Aşama~4'te
yapılacak son statik analiz, aynı kod tabanı üzerinde aynı komut zinciri
ile (\texttt{npm run analyze:stage2}) tekrarlanacaktır.

% ----------------------------------------------------------------------------
\begin{thebibliography}{00}
\bibitem{b1} SonarSource, ``SonarQube Documentation,''
  \url{https://docs.sonarsource.com/sonarqube/}, Erişim: 2026-05-07.
\bibitem{b2} ESLint, ``ESLint -- Pluggable JavaScript Linter,''
  \url{https://eslint.org/}, Erişim: 2026-05-07.
\bibitem{b3} G. A. Campbell, ``Cognitive Complexity: A New Way of
  Measuring Understandability,'' SonarSource S.A., Teknik Rapor, 2018.
\end{thebibliography}

\end{document}
